\شروع‌مسئله [10]

آرایه $A$ از اعداد طبیعی داده شده است. می دانیم هر عدد طبیعی حداکثر $1000$ بار در آرایه آمده است.

دو نوع کوئری داریم:
\شروع{شمارش}
\فقره
کوئری نوع اول این است که برای هر دو عدد طبیعی $n,r$ جایگاه $r$ امین $n$ در دنباله را بگوییم.

\فقره کوئری نوع دوم این است که عضوی مانند $i$ را در دنباله عوض کرده و $j$ را قرار دهیم.(تضمین می شود این تغییر شرط اولیه آرایه را نگه می دارد.)

\پایان{شمارش}
با استفاده از درهم‌سازی داده ساختاری طراحی کنید جواب هر پرسش را در $O(1)$ بدهد.  

\solution{
جواب مسئله یکتا نیست اما یکی از جواب هایی که این پیچیدگی زمانی را به ما می دهد داده ساختار زیر است:

از یک جدول درهم سازی زنحیره ای استفاده می کنیم و روی هر عضو لینک لیست های جدول یک پوینتر به یک $AVL\:TREE$ قرار می دهیم که در هر $AVL\:TREE$ جایگاه های عدد مربوط به آن را قرار می دهیم.

در این صورت برای هر پرسش نوع اول در $O(1)$ درخت راس پیدا می شود و چون هر درخت حداکثر $1000$ راس دارد در $O(1)$ می توانیم پاسخ پرسش را بدهیم.

به همین صورت چون اضافه کردن و حذف کردن از درخت ها در زمان $O(1)$ انجام می شود پرسش دوم هم در زمان مطلوب قابل انجام است.

در پاسخ زمان پیدا کردن یک درخت راس را $O(1)$ در نظر گرفتیم چون نسبت اندازه جدول و اندازه آرایه را میتوان عددی ثابت گرفت و بطور احتمالی طول هر لینک لیست ثابت است.

}